\medskip









\begin{enumerate}
\item 
	\begin{enumerate}
		\item 
Il y a 7 tirages contenant la lettre A, puis 

6 tirages contenant la lettre B (le tirage AB étant le même que le tirage BA), 

5 tirages contenant la lettre C, etc.

Il y a donc : $7 + 6 + 5 + 4 + 3 + 2 + 1 = \dfrac{7 \times 8}{2}= 7 \times 4 = 28$ tirages différents.
		\item 
Les tirages gagnant sont les 6 tirages contenant la lettre A et une consonne et les 6 contenanr la lettre E et une consonne : il y a donc $6 + 6 = 12$ tirages gagants.

La probabilité que le joueur gagne à ce jeu est donc égale à $\dfrac{12}{28} = \dfrac{4\times 3}{4 \times 7} = \dfrac{3}{7}$.
	\end{enumerate}
\end{enumerate}
	




\begin{enumerate}[resume]
\item 





	\begin{enumerate}
		\item 
On a $P(G = 10 - k) = \dfrac{3}{7}$ et $P(G = -k) = \dfrac{4}{7}$. D'où le tableau :

\begin{center}
$\begin{array}{|*{3}{c|}}\hline
G				&- k			&10 - k\\ \hline
P(G = ...)		&\dfrac{4}{7}	&\dfrac{3}{7}\rule[-3mm]{0mm}{9mm}\\ \hline
\end{array}$
\end{center}
		\item 
L'espérance mathématique de la variable aléatoire $G$ est égale à 
		
$E(G) = - k \times \dfrac{4}{7} + (10 - k) \times \dfrac{3}{7} = \dfrac{-4k + 30 - 3k}{7} = \dfrac{30 - 7k}{7}$.

Le jeu est favorable au joueur si :

$E(G) > 0 \iff \dfrac{30 - 7k}{7} \iff 30 - 7k > 0 \iff 7k < 30 \iff k < \dfrac{30}{7}$.

$\dfrac{30}{7} \approx 4,3$.

La somme payée au départ pour que le jeu reste favorable au joueur ne doit pas dépasser 4~\euro.	
	\end{enumerate}
\item 


	\begin{enumerate}
		\item 
Le tirage par un joueur est indépendant de celui des autres et chacun a une probabilité de gagner de $\dfrac{3}{7}$ ; $X$ suit donc une loi binomiale de paramètres $n = 10$ et $p = \dfrac{3}{7}$.
		\item 
On a $p(X = 4) = \binom{10}{4}\left(\dfrac{3}{7}\right)^4\left( 1 - \dfrac{3}{7}  \right)^{10-4} = 210 \times\left(\dfrac{3}{7}\right)^4 \times \left(\dfrac{4}{7}\right)^6 \approx \np{0,2466}$, soit 0,247 au millième près.
		\item  
La calculatrice donne $p(X\leqslant 4) \approx 0,560$ donc 

$p(X\geqslant 5) = 1 - p(X\leqslant 4) \approx (1 - 0,560)$.

Finalement : $p(X\geqslant 5) \approx 0,440$.
		
La probabilité qu'il y ait au moins 5 gagnants sur 10 joueurs est d'environ 0,440.
		\item 
La calculatrice donne :
		
$p(X \leqslant 5) \approx 0,782$ et $p(X \leqslant 6) \approx 0,921$, donc 
le plus petit entier $n$ tel que 

$P(X \leqslant n) \geqslant 0,9$ est donc $n = 6$.
	\end{enumerate}
\end{enumerate}


\medskip
